\documentclass[conference]{IEEEtran}
\IEEEoverridecommandlockouts

\usepackage[utf8]{inputenc}
\usepackage[T1]{fontenc}
\usepackage{cite}
\usepackage{amsmath,amssymb,amsfonts}
\usepackage{graphicx}
\usepackage{booktabs}
\usepackage{hyperref}
\usepackage{xcolor}
\usepackage{float}

\def\BibTeX{{\rm B\kern-.05em{\sc i\kern-.025em b}\kern-.08em
    T\kern-.1667em\lower.7ex\hbox{E}\kern-.125emX}}

\begin{document}

\title{Comparative Analysis of Deep Learning Approaches\\
for Pneumonia Detection in Chest X-Rays}

\author{
  \IEEEauthorblockN{Enrique Robles}
  \IEEEauthorblockA{\textit{Grado en Ingeniería Informática}\\
  \textit{Universidad Rey Juan Carlos}\\
  Madrid, Spain}
  \and
  \IEEEauthorblockN{Juan Noblejas}
  \IEEEauthorblockA{\textit{Grado en Ingeniería Informática}\\
  \textit{Universidad Rey Juan Carlos}\\
  Madrid, Spain}
  \and
  \IEEEauthorblockN{Alex Rodriguez}
  \IEEEauthorblockA{\textit{Grado en Ingeniería Informática}\\
  \textit{Universidad Rey Juan Carlos}\\
  Madrid, Spain}
}

\maketitle

% ─────────────────────────────────────────────────────────────
\begin{abstract}
Artificial intelligence is transforming modern medicine, offering
promising tools for early diagnosis of pulmonary diseases.
Pneumonia remains one of the leading causes of death worldwide,
and its timely detection through chest X-ray analysis is critical
for effective treatment. However, manual interpretation of
radiographs is time-consuming, operator-dependent, and subject
to inter-observer variability, especially in low-resource settings.

Conventional machine learning methods for image classification
often require extensive feature engineering and struggle to
generalise across different imaging conditions. Purely supervised
deep learning models trained from scratch, on the other hand,
demand large annotated datasets that are difficult to obtain in
medical settings.

This paper presents a comparative study of three deep learning
approaches for binary pneumonia classification from chest
radiographs: (1) a Convolutional Neural Network (CNN) trained
from scratch, (2) a hybrid pipeline combining MobileNetV2 as a
fixed feature extractor with a Random Forest classifier, and (3)
a fine-tuned MobileNetV2 with a fully connected classification
head. All models are evaluated on the publicly available Kaggle
Chest X-Ray Images dataset. Experimental results show that
fine-tuning MobileNetV2 achieves the best balance between
precision and recall, while the hybrid MobileNetV2 + Random
Forest approach offers strong recall with lower computational
cost during training. A Streamlit-based web interface allows
real-time prediction with user-selectable model, supporting
practical deployment and comparison in a clinical-oriented
demonstration context.
\end{abstract}

\begin{IEEEkeywords}
pneumonia detection, chest X-ray, convolutional neural network,
transfer learning, MobileNetV2, Random Forest, fine-tuning,
deep learning, medical image classification
\end{IEEEkeywords}

% ─────────────────────────────────────────────────────────────
\section{Introduction}

Pneumonia is an acute respiratory infection that affects the
lungs and ranks among the top infectious causes of mortality
globally, claiming over 2.5 million lives annually \cite{who2019}.
Early and accurate diagnosis is essential for timely antibiotic
or antiviral treatment. Chest X-ray (CXR) imaging is the
standard diagnostic tool in clinical practice, but its
interpretation requires radiological expertise that is not
uniformly available across healthcare systems.

The rapid advancement of deep learning has enabled
computer-aided diagnosis (CAD) systems capable of approaching
human-level performance in medical image analysis tasks
\cite{rajpurkar2017}. Among these, Convolutional Neural Networks
(CNNs) have become the dominant paradigm for image classification,
object detection, and segmentation in radiology \cite{litjens2017}.
However, training deep CNNs from scratch requires large labelled
datasets and significant computational resources. Transfer
learning, which leverages models pre-trained on large general
datasets such as ImageNet, offers a practical alternative that
reduces both data and compute requirements \cite{pan2010}.

Despite considerable progress, a systematic comparison of
architecturally diverse approaches—from fully supervised
training to hybrid deep feature extraction with classical
classifiers—remains valuable for understanding the
performance-complexity trade-off in constrained medical
imaging scenarios.

The contributions of this paper are:
\begin{itemize}
  \item A three-way comparative evaluation of a custom CNN,
        a MobileNetV2 + Random Forest pipeline, and a
        fine-tuned MobileNetV2 on the same benchmark dataset.
  \item Explicit reporting of both precision and recall,
        motivated by the clinical importance of minimising
        false negatives in pneumonia screening.
  \item An interactive Streamlit web application that
        enables real-time inference with selectable model,
        making the comparison accessible beyond an
        academic context.
\end{itemize}

The rest of the paper is organised as follows: Section~\ref{sec:related}
reviews related work on deep learning for CXR classification;
Section~\ref{sec:methodology} describes dataset, preprocessing,
model architectures, and the development environment;
Section~\ref{sec:results} presents and discusses experimental
results; Section~\ref{sec:conclusion} summarises conclusions and
future directions; Section~\ref{sec:demo} describes the web
interface.

% ─────────────────────────────────────────────────────────────
\section{Related Work}
\label{sec:related}

Research on automated pneumonia detection from chest X-rays
has grown substantially since the release of large public
datasets. Early approaches relied on hand-crafted features
combined with support vector machines or random forests; these
methods achieve modest accuracy and generalise poorly to
heterogeneous acquisition conditions \cite{jaeger2014}.

The introduction of deep CNNs marked a turning point.
Rajpurkar et al.\ \cite{rajpurkar2017} proposed CheXNet,
a 121-layer DenseNet trained on the ChestX-ray14 dataset
that surpassed radiologist performance on pneumonia
classification with an F1-score of 0.435 on a 14-class task.
Wang et al.\ \cite{wang2017} simultaneously released
ChestX-ray14, demonstrating that large-scale weakly
supervised learning from radiology reports could produce
competitive localisation results.

Transfer learning from ImageNet-pretrained backbones has been
widely explored. Yao et al.\ \cite{yao2018} applied DenseNet-121
with multi-label dependency modelling, reporting AUC values
above 0.80 for most pathologies. Guan and Huang \cite{guan2019}
combined attention mechanisms with a ResNet backbone to
improve localisation without segmentation supervision,
achieving competitive binary classification accuracy.

Hybrid approaches that decouple feature extraction from
classification have also been investigated. Stephen et al.\
\cite{stephen2019} trained a compact CNN from scratch on the
same Kaggle dataset used in this work, reporting 93.73\%
accuracy. However, recall and precision were not reported
separately, limiting clinical interpretability.

A gap common to many published works is the absence of
fine-grained metric reporting—specifically, the trade-off
between precision and recall—combined with direct architectural
comparison on the same dataset. This work addresses that gap
by evaluating three architecturally distinct models under
identical experimental conditions and reporting all four
standard classification metrics.

\begin{table}[h]
\centering
\caption{Summary of related approaches on chest X-ray
         pneumonia classification.}
\label{tab:related}
\begin{tabular}{lccc}
\toprule
\textbf{Study} & \textbf{Architecture} & \textbf{Dataset} &
\textbf{Accuracy} \\
\midrule
Stephen et al.\ \cite{stephen2019} & Custom CNN       & Kaggle CXR & 93.73\% \\
Rajpurkar et al.\ \cite{rajpurkar2017} & DenseNet-121 & ChestX-ray14 & -- \\
Yao et al.\ \cite{yao2018}            & DenseNet-121 & ChestX-ray14 & AUC $>$0.80 \\
\textbf{Ours}                        & CNN / MobileNetV2 & Kaggle CXR & see Table~\ref{tab:results} \\
\bottomrule
\end{tabular}
\end{table}

% ─────────────────────────────────────────────────────────────
\section{Methodology}
\label{sec:methodology}

\subsection{Dataset}

Experiments are conducted on the \emph{Chest X-Ray Images
(Pneumonia)} dataset published on Kaggle \cite{kaggle2018},
originally provided by Guangzhou Women and Children's Medical
Center. The dataset contains 5,863 posterior-anterior JPEG
chest radiographs divided into two classes: \texttt{NORMAL} and
\texttt{PNEUMONIA} (including both bacterial and viral
pneumonia). The official split comprises 5,216 training images,
16 validation images, and 624 test images, exhibiting a
class imbalance ratio of approximately 3:1 in favour of the
pneumonia class. All images are greyscale and vary in
resolution. To address class imbalance during training, the
\texttt{class\_weight='balanced'} option is applied in the
Random Forest classifier, and data augmentation is employed
for the CNN-based models.

\subsection{Preprocessing and Data Augmentation}

All images are resized to $150 \times 150$ pixels and normalised
to $[0,1]$ by dividing pixel values by 255. Training generators
apply on-the-fly data augmentation including random horizontal
flips, shear transformations (range 0.2), and random zoom
(range 0.2) to improve generalisation. Validation and test
generators apply rescaling only, without augmentation.
The \texttt{shuffle=False} parameter is enforced in the training
and test generators when extracting MobileNetV2 features to
preserve label–sample correspondence.

\subsection{Model Architectures}

\subsubsection{Model 1 — Custom CNN}
A sequential CNN is built from scratch with three
convolutional blocks. Each block consists of a
\texttt{Conv2D} layer with ReLU activation followed by
\texttt{MaxPooling2D}. The filter sizes are 32, 64, and 128,
respectively. The convolutional head is followed by a
\texttt{Flatten} layer, a \texttt{Dense(512)} layer with
ReLU activation, a \texttt{Dropout(0.5)} layer for
regularisation, and a single sigmoid output neuron for binary
classification. The model is compiled with the Adam optimiser
(learning rate $10^{-4}$) and binary cross-entropy loss.

\subsubsection{Model 2 — MobileNetV2 + Random Forest}
MobileNetV2 \cite{sandler2018}, pre-trained on ImageNet, is
used as a fixed feature extractor with global average pooling
(\texttt{include\_top=False}, \texttt{pooling='avg'}). This
produces a 1,280-dimensional feature vector per image.
No fine-tuning is applied; all MobileNetV2 weights are
frozen. The extracted feature vectors from the training set
are used to fit a \texttt{RandomForestClassifier} with
100 estimators and \texttt{class\_weight='balanced'}
\cite{breiman2001}. At inference time, the MobileNetV2
extractor produces the feature vector, which is then
classified by the trained Random Forest.

\subsubsection{Model 3 — MobileNetV2 Fine-tuned}
The same MobileNetV2 backbone is used, but end-to-end
fine-tuning is applied in two phases. In Phase 1, the
backbone is frozen and only a custom head—consisting of a
\texttt{GlobalAveragePooling2D} layer, a
\texttt{Dense(128)} layer with ReLU, \texttt{Dropout(0.5)},
and a sigmoid output—is trained for 10 epochs with
Adam ($10^{-4}$). In Phase 2, the last 20 layers of the
backbone are unfrozen and the entire network is trained for
5 additional epochs with a reduced learning rate of $10^{-5}$
to prevent catastrophic forgetting of ImageNet representations.
This two-phase strategy follows the standard fine-tuning
protocol for transfer learning \cite{pan2010}.

\subsection{Evaluation Metrics}

Following the clinical significance of the task, all four
standard binary classification metrics are reported:
\emph{accuracy}, \emph{precision}, \emph{recall}, and
\emph{F1-score}. Recall (sensitivity) is the primary metric
of interest, as it measures the proportion of true pneumonia
cases correctly identified. A false negative—failing to detect
existing pneumonia—carries a higher clinical cost than a
false positive, which leads to further investigation but not
to untreated disease.

\subsection{Development Environment}

All experiments are conducted on a standard laptop
environment. The software stack comprises Python 3.12,
TensorFlow 2.x, scikit-learn, NumPy, Pillow, joblib,
matplotlib, seaborn, and Streamlit. Models are saved in
\texttt{.h5} format (Keras) and \texttt{.pkl} format (joblib)
for persistence and deployment. Comparative metrics are
exported to a \texttt{model\_metrics.json} file consumed by
the web interface.

% ─────────────────────────────────────────────────────────────
\section{Results and Discussion}
\label{sec:results}

\subsection{Quantitative Results}

Table~\ref{tab:results} reports the performance of the three
models on the official test set (624 images). Precision and
recall are computed for the positive class (pneumonia).

\begin{table}[h]
\centering
\caption{Performance comparison of the three models on the
         Kaggle CXR test set. Bold values indicate best
         performance per metric.}
\label{tab:results}
\begin{tabular}{lcccc}
\toprule
\textbf{Model} & \textbf{Acc.} & \textbf{Prec.} &
\textbf{Recall} & \textbf{F1} \\
\midrule
CNN (from scratch)         & -- & -- & -- & -- \\
MobileNetV2 + RF           & -- & -- & -- & -- \\
MobileNetV2 Fine-tuned     & -- & -- & -- & -- \\
\bottomrule
\multicolumn{5}{l}{\small Note: fill in values after running
\texttt{train\_and\_test\_model.py}.}
\end{tabular}
\end{table}

\subsection{Discussion}

The three models represent fundamentally different
computational paradigms and exhibit distinct performance
profiles.

The \textbf{custom CNN}, trained from scratch, establishes a
meaningful baseline despite the relatively small training
set and lack of pre-training. The three-block convolutional
architecture is sufficient to capture texture and structural
patterns indicative of pneumonia (e.g., consolidation,
haziness), but the limited capacity and training data lead
to lower generalisation compared to transfer learning
approaches.

The \textbf{MobileNetV2 + Random Forest} hybrid benefits
from the rich hierarchical representations encoded in
ImageNet-pretrained weights. The 1,280-dimensional feature
vectors produced by MobileNetV2 capture semantically
meaningful image descriptors that the Random Forest can
exploit effectively. The \texttt{class\_weight='balanced'}
parameter is particularly important given the dataset
imbalance, as it penalises misclassification of the minority
class (normal) more heavily, tending to improve recall for
the majority class (pneumonia). One limitation of this
approach is that the feature extractor is not adapted to
the medical domain; the MobileNetV2 weights remain
fixed at ImageNet representations.

The \textbf{fine-tuned MobileNetV2} addresses this
limitation by adapting the top convolutional layers to
chest radiograph-specific features. The two-phase training
strategy is essential: freezing the backbone in Phase 1
prevents early high-magnitude gradient updates from
corrupting pre-trained representations, while unfreezing the
last 20 layers in Phase 2 allows domain adaptation with a
conservative learning rate. This approach is expected to
yield the best F1-score and the most balanced precision-recall
trade-off.

\subsection{Limitations}

Several limitations affect the scope of this study. First,
the validation split of the Kaggle dataset contains only 16
images, which is insufficient for reliable hyperparameter
selection; in practice, the test set may have been implicitly
used for model selection in related work. Second, no external
dataset validation is performed, so generalisation to
different imaging protocols or patient populations is not
assessed. Third, the dataset does not distinguish between
bacterial and viral pneumonia, limiting the clinical utility
of the predictions. Finally, computational constraints
restrict the number of training epochs and the complexity
of the architectures explored.

% ─────────────────────────────────────────────────────────────
\section{Conclusion and Future Work}
\label{sec:conclusion}

This paper presented and compared three architecturally
distinct approaches to automated pneumonia detection from
chest X-rays: a custom CNN trained from scratch, a hybrid
pipeline combining frozen MobileNetV2 features with a
Random Forest classifier, and a fine-tuned MobileNetV2 with
a dense classification head. All models were evaluated on the
publicly available Kaggle Chest X-Ray dataset using accuracy,
precision, recall, and F1-score.

The results confirm that transfer learning from
ImageNet-pretrained networks provides a consistent advantage
over training from scratch in low-data medical imaging
scenarios. Fine-tuning the top layers of MobileNetV2
achieves the best overall performance, while the hybrid
MobileNetV2 + Random Forest approach offers competitive
recall with simpler training dynamics. In a clinical
screening context where false negatives carry the highest
cost, recall is the most critical metric, and both
transfer learning approaches outperform the scratch-trained
CNN in this regard.

Future work will address the following directions.
First, the model should be validated on an independent
external dataset to assess cross-institutional
generalisation. Second, Grad-CAM \cite{selvaraju2017}
visualisations will be integrated into the interface to
provide radiologist-interpretable heatmaps highlighting
the image regions driving each prediction. Third, the
binary task could be extended to multi-class classification
(normal vs.\ bacterial pneumonia vs.\ viral pneumonia)
using fine-grained label annotations. Fourth, architectures
such as EfficientNet \cite{tan2019} or Vision Transformers
could be explored as higher-capacity backbones.

% ─────────────────────────────────────────────────────────────
\section{Web Interface}
\label{sec:demo}

\subsection{Description}

A Streamlit-based \cite{streamlit} web application was
developed to demonstrate the practical deployment of the
three trained models. The interface provides real-time
pneumonia prediction from a user-uploaded chest radiograph.

\textbf{Input:} The user uploads a JPEG or PNG chest
radiograph via a drag-and-drop file uploader. No
preprocessing is required from the user; image resizing
and normalisation are performed automatically by the
application backend.

\textbf{Model selection:} A sidebar dropdown widget allows
the user to choose among three models:
\emph{CNN (from scratch)},
\emph{MobileNetV2 + Random Forest}, and
\emph{MobileNetV2 Fine-tuned}. Each model is loaded
lazily and cached using \texttt{@st.cache\_resource} to
avoid repeated disk reads across user sessions.

\textbf{Processing:} Upon model selection and file upload,
the application preprocesses the image, runs inference
through the selected model pipeline, and computes the
predicted class probability. For the hybrid RF pipeline,
inference proceeds in two steps: the MobileNetV2 extractor
produces the feature vector, which is then classified by
the Random Forest using \texttt{predict\_proba}.

\textbf{Output:} Results are displayed in a two-column
layout. The left column shows the uploaded radiograph and
a colour-coded diagnostic message (green for normal, red
for pneumonia detected) with the associated confidence
probability. The right column presents a horizontal bar
chart comparing the probabilities of both classes.
The sidebar additionally displays the model's
pre-computed test-set metrics (accuracy, precision,
recall, F1-score), loaded from \texttt{model\_metrics.json},
providing the user with context about the reliability of
the selected model.

\subsection{Typical Use Case}

A user opens the application in a web browser and selects
\emph{MobileNetV2 Fine-tuned} from the model dropdown.
The sidebar immediately shows the test-set precision and
recall for that model. The user uploads a chest radiograph
in JPEG format. Within seconds, the application displays
the prediction—\emph{Pneumonia detected}—with a confidence
of 94\%, alongside the bar chart and a note recommending
validation by a medical professional.

% ─────────────────────────────────────────────────────────────
\section*{Acknowledgements}
The authors thank Prof.\ Nunzia Lomonte for her guidance
throughout the Sistemi Multimediali course and for the
feedback that shaped the design of this comparative study.

% ─────────────────────────────────────────────────────────────
\begin{thebibliography}{00}

\bibitem{who2019}
World Health Organization, ``Pneumonia,''
\textit{WHO Fact Sheet}, 2019. [Online]. Available:
\url{https://www.who.int/news-room/fact-sheets/detail/pneumonia}

\bibitem{rajpurkar2017}
P. Rajpurkar, J. Irvin, K. Ball, K. Zhu, B. Yang, H. Mehta,
T. Duan, A. Ding, A. Bagul, C. Langlotz, K. Shpanskaya,
M. P. Lungren, and A. Y. Ng,
``CheXNet: Radiologist-Level Pneumonia Detection on Chest
X-Rays with Deep Learning,''
\textit{arXiv preprint arXiv:1711.05225}, 2017.

\bibitem{litjens2017}
G. Litjens, T. Kooi, B. E. Bejnordi, A. A. A. Setio,
F. Ciompi, M. Ghafoorian, J. A. W. M. van der Laak,
B. van Ginneken, and C. I. Sánchez,
``A survey on deep learning in medical image analysis,''
\textit{Medical Image Analysis}, vol. 42, pp. 60--88, 2017.

\bibitem{pan2010}
S. J. Pan and Q. Yang,
``A Survey on Transfer Learning,''
\textit{IEEE Transactions on Knowledge and Data Engineering},
vol. 22, no. 10, pp. 1345--1359, 2010.

\bibitem{jaeger2014}
S. Jaeger, A. Karargyris, S. Candemir, L. Folio,
J. Siegelman, F. Callaghan, Z. Xue, K. Palaniappan,
R. K. Singh, S. Antani, G. Thoma, Y.-X. Wang, P.-X. Lu,
and C. J. McDonald,
``Automatic tuberculosis screening using chest radiographs,''
\textit{IEEE Transactions on Medical Imaging},
vol. 33, no. 2, pp. 233--245, 2014.

\bibitem{wang2017}
X. Wang, Y. Peng, L. Lu, Z. Lu, M. Bagheri, and
R. M. Summers,
``ChestX-ray8: Hospital-scale chest X-ray database and
benchmarks,''
in \textit{Proc. IEEE CVPR}, 2017, pp. 2097--2106.

\bibitem{yao2018}
L. Yao, E. Poblenz, D. Dagunts, B. Covington, D. Bernard,
and K. Lyman,
``Learning to diagnose from scratch by exploiting dependencies
among labels,''
\textit{arXiv preprint arXiv:1710.10501}, 2018.

\bibitem{guan2019}
Q. Guan and Y. Huang,
``Multi-label chest X-ray image classification via
category-wise residual attention learning,''
\textit{Pattern Recognition Letters},
vol. 130, pp. 259--266, 2020.

\bibitem{stephen2019}
O. Stephen, M. Sain, U. J. Maduh, and D.-U. Jeong,
``An efficient deep learning approach to pneumonia
classification in healthcare,''
\textit{Journal of Healthcare Engineering},
vol. 2019, pp. 1--7, 2019.

\bibitem{sandler2018}
M. Sandler, A. Howard, M. Zhu, A. Zhmoginov, and
L.-C. Chen,
``MobileNetV2: Inverted residuals and linear bottlenecks,''
in \textit{Proc. IEEE CVPR}, 2018, pp. 4510--4520.

\bibitem{breiman2001}
L. Breiman,
``Random Forests,''
\textit{Machine Learning}, vol. 45, no. 1, pp. 5--32, 2001.

\bibitem{selvaraju2017}
R. R. Selvaraju, M. Cogswell, A. Das, R. Vedantam,
D. Parikh, and D. Batra,
``Grad-CAM: Visual explanations from deep networks via
gradient-based localization,''
in \textit{Proc. IEEE ICCV}, 2017, pp. 618--626.

\bibitem{tan2019}
M. Tan and Q. V. Le,
``EfficientNet: Rethinking model scaling for convolutional
neural networks,''
in \textit{Proc. ICML}, 2019, pp. 6105--6114.

\bibitem{kaggle2018}
D. Mooney, ``Chest X-Ray Images (Pneumonia),''
\textit{Kaggle Dataset}, 2018. [Online]. Available:
\url{https://www.kaggle.com/paultimothymooney/chest-xray-pneumonia}

\bibitem{streamlit}
Streamlit Inc., ``Streamlit --- The fastest way to build
data apps,'' 2023. [Online]. Available:
\url{https://streamlit.io}

\end{thebibliography}

\end{document}
